\documentclass[conference]{IEEEtran}

% Core packages
\usepackage{cite}
\usepackage{amsmath,amssymb,amsfonts}
\usepackage{graphicx}
\usepackage{textcomp}
\usepackage{xcolor}
\usepackage{booktabs}
\usepackage{url}

\begin{document}

\title{A Data-Driven Framework for Tourism Review Analytics in Sri Lanka}

\author{%
\IEEEauthorblockN{First Author}
\IEEEauthorblockA{%
Department of Computer Science\\
University Name\\
City, Country\\
email@domain.com}
\and
\IEEEauthorblockN{Second Author}
\IEEEauthorblockA{%
Department of Computer Science\\
University Name\\
City, Country\\
email@domain.com}
\and
\IEEEauthorblockN{Third Author}
\IEEEauthorblockA{%
Department of Computer Science\\
University Name\\
City, Country\\
email@domain.com}
}

\maketitle

\begin{abstract}
This paper presents a reproducible end-to-end framework for tourism review analysis using structured and unstructured user-generated content. The workflow includes dataset preparation, location and temporal feature engineering, sentiment-oriented labeling, statistical analysis, and model validation. We process review text and metadata to construct robust predictors such as review delay, review length, user contribution behavior, and location-derived signals. Experimental results on a real-world tourism review dataset show that the proposed feature set improves predictive performance and supports actionable insights for destination management. The framework is designed for reproducibility and can be adapted to other regional tourism datasets with minimal changes.
\end{abstract}

\begin{IEEEkeywords}
Tourism analytics, sentiment analysis, feature engineering, machine learning, reproducibility
\end{IEEEkeywords}

\section{Introduction}
Tourism platforms generate large volumes of user reviews that contain valuable behavioral and sentiment signals. Transforming these raw signals into reliable analytical outputs requires careful data preparation, domain-aware feature engineering, and rigorous validation.

This study develops a practical framework for tourism review analytics with emphasis on reproducibility and conference-grade reporting. The framework is demonstrated on Sri Lankan tourism review data and supports both statistical interpretation and predictive modeling.

\textbf{Contributions of this paper are as follows:}
\begin{itemize}
    \item A reproducible preprocessing pipeline for tourism reviews with deterministic handling of location and country extraction.
    \item A comprehensive feature engineering strategy combining location, temporal, and text-derived indicators.
    \item An empirical evaluation protocol for comparative model validation and interpretation.
\end{itemize}

\section{Related Work}
Prior work in tourism analytics has used sentiment mining, aspect extraction, and predictive models on online reviews~\cite{liu2012sentiment,cambria2017sentic}. Recent studies emphasize combining metadata with textual features to improve model robustness~\cite{zhang2021review}. However, many pipelines are difficult to reproduce due to inconsistent preprocessing and weak reporting of feature construction.

Our work addresses this gap through a transparent feature engineering workflow and explicit export of processed datasets for downstream tasks.

\section{Methodology}
\subsection{Dataset and Preprocessing}
The dataset consists of tourism reviews with fields such as location name, city, user locale, rating, review text, and timestamps. Records are cleaned for missing values, duplicates, and type inconsistencies.

Location normalization includes province and district extraction from raw location strings using curated mappings and rule-based parsing. User country inference combines direct string matching, locale-based fallback, and cache-backed resolution to ensure deterministic outputs.

\subsection{Feature Engineering}
The following feature groups are generated:
\begin{itemize}
    \item \textbf{Location features:} province, district, city, and location type.
    \item \textbf{Temporal features:} travel and publication month/year; review delay in days.
    \item \textbf{Behavioral features:} user contributions and helpful votes.
    \item \textbf{Text features:} title length and review length.
    \item \textbf{Label feature:} rating class (Negative, Neutral, Positive).
\end{itemize}

\subsection{Modeling and Validation}
A supervised learning setup is used for rating class prediction. Data is split into training and testing partitions, and cross-validation is applied for robust comparison. Performance is reported using Accuracy, Precision, Recall, and F1-score.

\section{Experimental Setup}
\subsection{Implementation Details}
Experiments are implemented in Python with common data science libraries. The preprocessing and feature generation pipeline is executed before model training to avoid leakage.

\subsection{Evaluation Metrics}
For classification, we report:
\begin{equation}
\text{Accuracy} = \frac{TP + TN}{TP + TN + FP + FN}
\end{equation}
where $TP$, $TN$, $FP$, and $FN$ denote true positives, true negatives, false positives, and false negatives.

\section{Results and Discussion}
Table~\ref{tab:results} provides a template for model comparison.

\begin{table}[htbp]
\caption{Model Performance Comparison}
\label{tab:results}
\centering
\begin{tabular}{lcccc}
\toprule
Model & Accuracy & Precision & Recall & F1-score \\
\midrule
Baseline Model & 0.000 & 0.000 & 0.000 & 0.000 \\
Model A & 0.000 & 0.000 & 0.000 & 0.000 \\
Model B & 0.000 & 0.000 & 0.000 & 0.000 \\
\bottomrule
\end{tabular}
\end{table}

The final paper should replace placeholder values with actual experimental results and include error analysis and ablation discussion.

\section{Threats to Validity}
Potential threats include dataset representativeness, annotation assumptions in rating-class conversion, and location parsing noise in ambiguous records. External validity should be assessed by testing the framework on additional tourism datasets.

\section{Conclusion and Future Work}
This paper introduced a reproducible framework for tourism review analytics integrating preprocessing, feature engineering, and model validation. Future work includes multilingual sentiment modeling, aspect-level opinion mining, and temporal drift analysis.

\section*{Acknowledgment}
This research was conducted as part of undergraduate research in data science. Replace this section with funding and institutional acknowledgments as applicable.

\bibliographystyle{IEEEtran}
\bibliography{references}

\end{document}
